\section*{Abstract}\label{sec:abstract}\justifying{}
{\color{red}
\textbf{Background/Objectives:} Dapagliflozin is an SGLT2 inhibitor prescribed for the management of type 2 diabetes mellitus. The drug lowers blood glucose levels by increasing urinary glucose excretion (UGE). Despite established efficacy, dapagliflozin demonstrates significant inter-individual variability in pharmacokinetics (PK) and pharmacodynamics (PD), with potential impact on treatment outcomes. \textbf{Methods:} To evaluate the sources of variability and to support patient stratification and model-informed individualized therapy, we developed a physiologically based pharmacokinetic/pharmacodynamic (PBPK/PD) model of dapagliflozin using curated data from 28 clinical studies. This framework integrates absorption, distribution, metabolism, excretion, and pharmacodynamics, and accounts for key determinants of variability including renal and hepatic function, and food effects. \textbf{Results:} The simulations reproduced dose-dependent pharmacokinetics with predicted Cmax and AUC values typically within 10--15\% of observed data. Renal impairment reduced UGE by 40--60\% despite modest changes in plasma exposure, while hepatic impairment produced only small shifts in PK and PD. The model also reproduced the fed-state reduction of peak concentrations, consistent with the 30--50\% decrease reported clinically. \textbf{Conclusions:} All model files, code, and curated datasets are openly available in line with FAIR standards and Open Science practices, enabling transparent and reproducible analyses and providing a mechanistic basis for individualized therapy in type 2 diabetes.
Enalapril is a medication used to treat high blood pressure and a num ber of other cardiovascular conditions. The renin-angiotensin-aldosterone system (RAAS) is the main component involved in the regulation of blood pressure in the body. One of the key steps in this process is the conversion of angiotensin I to angiotensin II by angiotensin converting enzyme(ACE). Enalapril is an ACE inhibitor and therefore helps to lower blood pressure. Pharmacokinetics is the branch of pharmacology that studies how the body absorbs, distributes, metabolizes, and eliminates drugs over time. After administration as enalapril maleate, it is absorbed in the intestine and converted by carboxylesterase 1 (CES1) in the liver to the active metabolite, enalaprilat. Enalapril and enalaprilat are eliminated from the body by renal clearance. A physiologically based pharmacokinetic (PBPK) model was developed to investigate the pharmacokinetics of enalapril and the factors influencing it. As part of the work, an extensive database of enalapril pharmacokinetics consisting of data from 49 clinical trials was established and used to parameterize and validate the computational model. The model was used to investigate the effect of renal impairment, hepatic impairment and changes in CES1 activity on the pharmacokinetics of enalapril and enalaprilat, as these three factors are known to have a major influence on the pharmacokinetics of enalapril. The model shows good agreement with a wide range of enalapril and enalaprilat data in serum, plasma and urine under different conditions (healthy, renal impairment, hepatic impairment, CES1 mutations). The model is available in SBML under a CC-BY 4.0 license with all data freely available from the PK-DB pharmacokinetics database.}

Keywords: hypertension, enalapril, pharmacokinetics, ACE inhibitor, renal impairment, hepatic impairment, CES1